\rtmtight
\begin{tblr}{width=\textwidth,colspec={|Q[c,m,2.1cm]|X[l,m]|},hlines,vlines,hspan=minimal,rowsep=1.5pt,colsep=3pt,row{1}={bg=headgray,font=\bfseries}}
\SetCell{c,m}\hfil\logo\hfil & \textbf{POLITEKNIK NEGERI MALANG}\par\textbf{JURUSAN TEKNOLOGI INFORMASI}\par\textbf{PROGRAM STUDI: D4 TEKNIK INFORMATIKA} \\
\SetCell[c=2]{l} \textbf{RENCANA TUGAS MAHASISWA (RTM) DAN RUBRIK PENILAIAN} & \\
Nama Mata Kuliah & Critical Thinking dan Problem Solving \\
Kode Mata Kuliah & RTI251003 \quad \textbf{sks / Jam perminggu:} 2 SKS/4 jam \quad \textbf{Semester:} 1 \\
Dosen Pengampu & \begin{enumerate}\item Vivi Nur Wijayaningrum, S.Kom., M.Kom.\item Elok Nur Hamdana, S.T., M.T.\item Rokhimatul Wakhidah, S.Pd., M.T.\item Mungki Astiningrum, S.T., M.Kom.\item Bagas Satya Dian Nugraha, S.T., M.T.\end{enumerate} \\
\end{tblr}
\assessmentblock{Tugas Analisis Argumen dan Pemecahan Masalah}{Tugas terstruktur mingguan}{SCPMK0401-00301, SCPMK0401-00302, SCPMK0401-00303}{Mahasiswa mengerjakan rangkaian Tugas 1--10 berupa analisis argumen, klasifikasi fakta/opini, evaluasi bias, penerapan teknik pemecahan masalah, dan perumusan solusi ilmiah dari studi kasus.}{Menganalisis studi kasus, mengklasifikasikan pernyataan, mengevaluasi argumen, menerapkan teknik problem solving (root cause, brainstorming, decision trees), dan menuliskan solusi berbasis bukti.}{Lembar jawaban tugas, hasil analisis studi kasus, dan draf tulisan solusi ilmiah.}{Ketepatan analisis argumen, fakta vs opini, bias, dan kesalahan logika & 10 & Premis, kesimpulan, fakta/opini, bias, dan kesalahan logika diidentifikasi tepat pada setiap kasus disertai alasan yang runtut. & Sebagian besar unsur argumen teridentifikasi benar, tetapi alasan atau identifikasi bias masih dangkal atau tidak konsisten. & Klasifikasi fakta/opini banyak keliru atau analisis argumen tidak menunjukkan pemahaman struktur premis-kesimpulan. \\ \\
Ketepatan penerapan teknik pemecahan masalah & 8 & Teknik (5-Why's, fishbone, problem tree, brainstorming, decision trees) dipilih sesuai kasus dan diterapkan dengan langkah lengkap hingga akar masalah dan alternatif solusi. & Teknik diterapkan pada kasus, tetapi pemilihan teknik kurang tepat atau langkah analisis belum tuntas sampai akar masalah. & Teknik tidak sesuai kasus, langkahnya tidak runtut, atau hasil analisis tidak terkait masalah. \\ \\
Perumusan solusi ilmiah berbasis data & 2 & Solusi dirumuskan logis, didukung data/bukti relevan, dan dituliskan dalam draf yang terstruktur. & Solusi cukup relevan, tetapi dukungan data atau struktur penulisan masih lemah. & Solusi asumtif tanpa dukungan data atau tidak menjawab masalah yang dirumuskan. \\}

\assessmentblock{Kuis Konsep Berpikir Kritis}{Kuis (ujian teori pilihan ganda, closed book)}{SCPMK0401-00301}{Mahasiswa mengerjakan kuis pilihan ganda yang mencakup materi pekan 1--4: konsep CTPS, fakta dan opini, struktur argumen, serta strategi pemecahan masalah dasar.}{Mengerjakan soal pilihan ganda secara mandiri dalam waktu yang ditentukan dengan sifat closed book.}{Lembar jawaban kuis pilihan ganda.}{Penguasaan konsep berpikir kritis dan CTPS & 6 & Menjawab benar soal konsep, karakteristik pemikir kritis, dan peran CTPS dengan tingkat ketepatan tinggi (di atas 80\% butir terkait). & Menjawab benar sebagian besar soal konsep (60--80\% butir terkait). & Jawaban benar kurang dari 60\% butir konsep, menunjukkan pemahaman yang belum memadai. \\ \\
Ketepatan identifikasi fakta, opini, dan klaim & 5 & Mengklasifikasikan pernyataan sebagai fakta, opini, atau klaim dengan tepat pada hampir semua butir soal. & Klasifikasi sebagian besar benar, tetapi masih tertukar pada pernyataan yang ambigu. & Klasifikasi banyak keliru pada butir fakta/opini/klaim. \\ \\
Ketepatan mengenali bias dan kesalahan logika & 4 & Mengenali jenis bias dan kesalahan penalaran pada soal kasus dengan tepat dan konsisten. & Mengenali sebagian bias/kesalahan logika, tetapi masih keliru pada kasus yang mirip. & Tidak dapat membedakan penalaran valid dan tidak valid pada butir soal. \\}

\assessmentblock{UTS Evaluasi Argumen dan Problem Solving}{UTS (ujian teori pilihan ganda, closed book)}{SCPMK0401-00301, SCPMK0401-00302}{Mahasiswa mengerjakan UTS yang mencakup materi pekan 1--7: analisis dan evaluasi argumen serta penerapan teknik pemecahan masalah pada kasus.}{Mengerjakan soal pilihan ganda secara mandiri dalam waktu yang ditentukan dengan sifat closed book.}{Lembar jawaban UTS pilihan ganda.}{Ketepatan evaluasi argumen dan validitas penalaran & 5 & Menilai validitas argumen, mengidentifikasi premis-kesimpulan, dan mengenali kesalahan penalaran pada soal kasus dengan benar dan konsisten. & Sebagian besar evaluasi argumen benar, tetapi masih keliru pada kasus penalaran yang kompleks. & Evaluasi argumen banyak keliru atau menilai validitas tanpa dasar logika yang benar. \\ \\
Ketepatan penerapan teknik problem solving pada kasus & 5 & Memilih dan menerapkan teknik pemecahan masalah (inferensi, reasoning matematis/grafis, decision trees) yang tepat untuk setiap soal kasus. & Teknik yang dipilih umumnya sesuai, tetapi penerapannya pada sebagian kasus belum tepat. & Pemilihan atau penerapan teknik tidak sesuai kasus sehingga jawaban banyak keliru. \\}

\assessmentblock{PBL Draf Dokumen Ilmiah Solusi Permasalahan}{Project Based Learning}{SCPMK0401-00303}{Mahasiswa secara berkelompok menyusun draf dokumen ilmiah (pekan 13--16) yang memuat latar belakang, rumusan masalah, analisis data, solusi, dan penyempurnaan berdasarkan feedback.}{Menyusun struktur proposal, membangun argumentasi ilmiah berbasis data, memadukan logika antarbagian dokumen, dan menyempurnakan solusi berdasarkan evaluasi kritis dan feedback.}{Draf dokumen ilmiah solusi permasalahan beserta catatan revisi.}{Struktur proposal dan perumusan masalah & 7 & Latar belakang, rumusan masalah, tujuan, dan sistematika dokumen lengkap, runtut, dan saling terhubung. & Struktur dokumen lengkap, tetapi keterkaitan antara latar belakang, rumusan masalah, dan tujuan belum kuat. & Struktur dokumen tidak lengkap atau rumusan masalah tidak jelas. \\ \\
Kualitas argumentasi ilmiah berbasis data & 8 & Setiap klaim solusi didukung data/bukti relevan dengan logika yang terpadu dari masalah hingga solusi. & Argumentasi cukup logis, tetapi sebagian klaim belum didukung data atau terdapat lompatan logika. & Argumentasi tidak didukung data atau solusi tidak terkait dengan rumusan masalah. \\ \\
Penyempurnaan solusi berdasarkan evaluasi kritis dan feedback & 5 & Feedback dianalisis, revisi diprioritaskan, dan perbaikan terlihat jelas pada draf akhir disertai catatan revisi. & Sebagian feedback ditindaklanjuti, tetapi dokumentasi revisi atau kedalaman perbaikan masih terbatas. & Tidak ada tindak lanjut bermakna terhadap feedback atau draf akhir tidak berubah. \\}

\assessmentblock{UAS Presentasi Akhir Solusi}{UAS (ujian presentasi kelompok)}{SCPMK0401-00303}{Mahasiswa mempresentasikan hasil analisis permasalahan dan solusi yang dikembangkan secara sistematis, serta menjawab pertanyaan dan memberikan klarifikasi berbasis bukti.}{Menyiapkan bahan presentasi dari dokumen ilmiah, memaparkan latar belakang, rumusan masalah, analisis data, dan solusi, kemudian menjawab pertanyaan penguji dengan argumentasi berbasis data.}{Bahan presentasi (slide), dokumen ilmiah final, dan sesi tanya jawab.}{Kejelasan struktur penyajian solusi & 12 & Presentasi runtut dari latar belakang, rumusan masalah, analisis, hingga solusi dengan alur yang mudah diikuti dan tepat waktu. & Struktur presentasi lengkap, tetapi urutan atau proporsi antarbagian kurang seimbang. & Presentasi melompat-lompat, bagian penting hilang, atau tidak menjelaskan rumusan masalah dan solusi. \\ \\
Kualitas argumentasi dan keterkaitan dengan data & 13 & Setiap gagasan solusi dijelaskan dengan argumen kohesif yang didukung data/fakta dari analisis yang dilakukan. & Argumentasi cukup jelas, tetapi sebagian gagasan belum didukung data atau keterkaitannya lemah. & Gagasan tidak didukung data atau argumen tidak konsisten dengan hasil analisis. \\ \\
Kemampuan menjawab pertanyaan dan klarifikasi berbasis bukti & 10 & Menjawab pertanyaan dengan tepat, tenang, dan merujuk pada bukti/data dalam dokumen ilmiah. & Sebagian pertanyaan dijawab tepat, tetapi jawaban lain kurang didukung bukti. & Tidak mampu menjawab pertanyaan atau jawaban tidak terkait bukti dan materi presentasi. \\}

\begin{tblr}{width=\textwidth,colspec={|Q[l,m,3.2cm]|X[l,m]|},hlines,vlines,hspan=minimal,row{1-3}={bg=headgray},rowsep=1.5pt,colsep=3pt}
Jadwal Pelaksanaan: & Minggu 1--17 sesuai RPS. Setiap evaluasi memiliki RTM dan rubrik multi-kriteria. \\
Lain-Lain Yang Diperlukan: & LMS, studi kasus, dokumen ilmiah, perangkat lunak office, dan perangkat presentasi. \\
Pustaka: & Mengacu pustaka utama dan pendukung pada bagian RPS. \\
\end{tblr}
